\documentclass[12pt,a4paper]{article}
\usepackage[utf8]{inputenc}
\usepackage[spanish]{babel}
\usepackage{geometry}
\usepackage{xcolor}
\usepackage{titlesec}
\usepackage{fancyhdr}
\usepackage{enumitem}
\usepackage{tikz}
\usetikzlibrary{shapes.geometric, arrows.meta, positioning, shadows}

% Configuración de márgenes e identidad visual
\geometry{left=2.54cm,right=2.54cm,top=2.54cm,bottom=2.54cm}
\definecolor{rojoUMB}{HTML}{8B0000}
\definecolor{grisFondo}{HTML}{F8F9FA}

% Configuración de encabezados
\pagestyle{fancy}
\fancyhf{}
\lhead{\footnotesize Manual de Usuario - PWA BioGALF}
\rhead{\footnotesize Fisioterapia Neurológica}
\lfoot{\footnotesize © 2026 Plataforma de Apoyo a la Decisión Clínica}
\rfoot{\thepage}

% Estilo de títulos
\titleformat{\section}{\color{rojoUMB}\normalfont\Large\bfseries}{\thesection}{1em}{}
\titleformat{\subsection}{\color{rojoUMB}\normalfont\large\bfseries}{\thesubsection}{1em}{}

\title{
    \vspace{-2cm}
    \textcolor{rojoUMB}{\textbf{MANUAL DE USUARIO ESTRUCTURADO}}\\
    \large \textit{Sistema BioGALF para Apoyo a la Decisión Clínica}
}
\author{}
\date{1 de abril de 2026}

\begin{document}

\maketitle
\vspace{-1.5cm}

\section{1. Flujo de Trabajo Secuencial}
Para garantizar la integridad de los datos, la aplicación está diseñada con un flujo lógico estricto.

\vspace{0.5cm}
\begin{center}
\begin{tikzpicture}[
    node distance=2cm,
    auto,
    block/.style={rectangle, draw=rojoUMB, fill=grisFondo, text width=4cm, text centered, rounded corners, minimum height=1.2cm, drop shadow, font=\footnotesize},
    line/.style={draw, -{Latex[length=2.5mm]}, rojoUMB, thick}
]
    % Nodos
    \node [block] (acceso) {\textbf{1. Escudo Profesional}\\Aceptación legal};
    \node [block, below=0.8cm of acceso] (perfil) {\textbf{2. Perfil Paciente}\\Datos y entorno};
    \node [block, below=0.8cm of perfil] (bateria) {\textbf{3. Batería Clínica}\\TUG, BBS, FGA};
    \node [block, right=1.5cm of bateria] (cueing) {\textbf{4. Metrónomo}\\Estimulación rítmica};
    \node [block, below=0.8cm of bateria] (guardar) {\textbf{5. Guardar en Nube}\\Registro validado};

    % Líneas
    \path [line] (acceso) -- (perfil);
    \path [line] (perfil) -- (bateria);
    \path [line] (bateria) -- (guardar);
    \path [line, dashed] (bateria) -- (cueing);
    \path [line, dashed] (cueing) |- (guardar);
\end{tikzpicture}
\end{center}
\vspace{0.5cm}

\begin{enumerate}
    \item \textbf{Aceptar acceso profesional:} Lea el aviso legal y acepte para habilitar la Guía de Intervención.
    \item \textbf{Llenar perfil del paciente:} Registre el nombre, el estadio de Parkinson y la evaluación socio-ambiental antes de iniciar las pruebas.
    \item \textbf{Ejecutar batería clínica:} Aplique TUG Cognitivo, Berg, FGA o UPDRS.
    \item \textbf{Usar Cueing con Metrónomo:} Active el estímulo rítmico según el objetivo terapéutico.
    \item \textbf{Guardar valoración:} Finalice guardando la sesión (se habilita tras 2 pruebas).
\end{enumerate}

\section{2. Explicación de Herramientas Clínicas}
\textit{Nota: La plataforma automatiza todos los cálculos y rangos de riesgo. No use calculadoras externas.}

\subsection*{TUG Cognitivo (Cronómetro Integrado)}
Presione \textbf{Start} al inicio y \textbf{Stop} al finalizar. El sistema evalúa automáticamente el punto de corte clínico (12.6 segundos) para detectar el riesgo de caída.

\subsection*{Escala de Berg (BBS) y FGA}
Seleccione los valores con los botones rápidos (0 al 4). El puntaje total se suma en tiempo real y la interpretación de riesgo se genera automáticamente.

\section{3. Nueva Herramienta de Cueing (Metrónomo)}
Diseñada para facilitar la regularidad de la marcha y combatir el \textit{Freezing of Gait} (congelamiento).

\vspace{0.5cm}
\begin{center}
\begin{tikzpicture}
    % Celular simulado
    \draw[rounded corners=4mm, thick, rojoUMB, fill=grisFondo] (0,0) rectangle (6,4);
    
    % Elementos del metrónomo
    \node[font=\bfseries, rojoUMB] at (3,3.3) {Estimulación Sensorial};
    \draw[thick, rojoUMB] (1,2.5) -- (5,2.5); % Slider
    \filldraw[rojoUMB] (2.5,2.5) circle (3pt); % Handle
    \node[font=\footnotesize] at (3,2.1) {Ajuste: 40 - 120 BPM};
    
    % Botón
    \draw[rounded corners, fill=rojoUMB!80, draw=none] (2,0.8) rectangle (4,1.6);
    \node[text=white, font=\footnotesize\bfseries] at (3,1.2) {INICIAR};
    
    % Ondas de sonido y luz
    \draw[rojoUMB, thick] (6.5,2.5) arc (-45:45:0.5);
    \draw[rojoUMB, thick] (6.8,2.8) arc (-45:45:0.8);
    \node[font=\footnotesize, text width=3cm] at (8.5,2) {\textbf{Pulso Sincronizado}\\- Beep Auditivo\\- Parpadeo Visual};
\end{tikzpicture}
\end{center}
\vspace{0.5cm}

\begin{itemize}
    \item \textbf{Ajuste BPM:} Use el control deslizante (40 a 120 BPM). Comience en un valor confortable (ej. 60) y ajuste según la estabilidad del paciente.
    \item \textbf{Señal Auditiva (Beep):} Marca el pulso temporal para sincronizar los pasos.
    \item \textbf{Señal Visual (Parpadeo):} Refuerza el ritmo para pacientes con fatiga atencional.
\end{itemize}

\section{4. Condición de Guardado Estricto}
El botón \textbf{Guardar Valoración} permanece deshabilitado por seguridad de datos. \textbf{Se habilitará únicamente tras completar al menos 2 pruebas de la batería clínica.} Esta regla protege la integridad del registro en la base de datos y evita historiales incompletos.

\end{document}